\section*{Chapter \ref{ch:g8:growth-and-compounding} --- Growth and Compounding}

\begin{solution}{exo:g8:growth-and-compounding:1}
South Korea: $1\,038$ dollars in 1960 and $37\,048$ in 2024. Argentina: $7\,397$ and $12\,774$.
\end{solution}

\begin{solution}{exo:g8:growth-and-compounding:2}
$(20\,600 - 20\,000) \div 20\,000 = 0.03$: $3\,\%$.
\end{solution}

\begin{solution}{exo:g8:growth-and-compounding:3}
$1\,050$; $1\,050 \times 1.05 = 1\,102.50$; $1\,102.50 \times 1.05 \approx 1\,157.63$.
\end{solution}

\begin{solution}{exo:g8:growth-and-compounding:4}
$70 \div 7 = 10$ years.
\end{solution}

\begin{solution}{exo:g8:growth-and-compounding:5}
$1.02^{35} \approx 2.00$: at $2\,\%$ a year, a quantity doubles in 35 years, as the rule of 70 says.
\end{solution}

\begin{solution}{exo:g8:growth-and-compounding:6}
$10\,000 \times 1.01^{20} \approx 12\,202$ dalens; $10\,000 \times 1.03^{20} \approx 18\,061$ dalens.
\end{solution}

\begin{solution}{exo:g8:growth-and-compounding:7}
$1.0575^{64} \approx 35.8$, close to $35.7$ (the rate $5.75\,\%$ is rounded).
\end{solution}

\begin{solution}{exo:g8:growth-and-compounding:8}
(a) \omterm{def:g8:growth-and-compounding:doubling-time}{Doubling time} $70 \div 5 = 14$ years; 28 years is two doublings: about $2\,000 \times 2 \times 2 =
8\,000$ dollars. (b) $2\,000 \times 1.05^{28} \approx 2\,000 \times 3.920 \approx 7\,840$ dollars.
\end{solution}

\begin{solution}{exo:g8:growth-and-compounding:9}
$(8\,048 - 11\,954) \div 11\,954 = -3\,906 \div 11\,954 \approx -0.327$: a fall of $32.7\,\%$. It does
not fit: the model predicts growth every year at the same rate.
\end{solution}

\begin{solution}{exo:g8:growth-and-compounding:10}
\begin{center}
\begin{tabular}{lrrrrr}
\toprule
Years & 0 & 10 & 20 & 30 & 40 \\
\midrule
Poorer ($6\,\%$) & $5\,000$ & $8\,954$ & $16\,036$ & $28\,717$ & $51\,429$ \\
Richer ($2\,\%$) & $20\,000$ & $24\,380$ & $29\,719$ & $36\,227$ & $44\,160$ \\
\bottomrule
\end{tabular}
\end{center}
The crossing is between 30 and 40 years. After 36 years: $40\,736$ against $40\,798$; after 37 years:
$43\,180$ against $41\,614$. The poorer country passes the richer one after 37 years.
\end{solution}

\begin{solution}{exo:g8:growth-and-compounding:11}
It doubles: from $70 \div 4 = 17.5$ years to $70 \div 2 = 35$ years.
\end{solution}

\begin{solution}{exo:g8:growth-and-compounding:12}
The rule only turns a \omterm{def:g8:growth-and-compounding:growth-rate}{growth rate} into a \omterm{def:g8:growth-and-compounding:doubling-time}{doubling time}. It says nothing about what produced the rate:
more \omterm{def:g8:production-and-productivity:capital-per-worker}{capital per worker}, better methods, more schooling, the rules under which people trade and
invest, or the luck of a period. Those causes are studied in later chapters and volumes.
\end{solution}

\begin{solution}{pb:g8:growth-and-compounding:1}
\textbf{1.} 2000: $24\,680 \div 10\,400 \approx 2.4$ times. 2025: $31\,470 \div 27\,700 \approx 1.1$
times.
\textbf{2.} Brenland: $(31\,470 - 24\,680) \div 24\,680 = 6\,790 \div 24\,680 \approx 27.5\,\%$.
Kestria: $(27\,700 - 10\,400) \div 10\,400 = 17\,300 \div 10\,400 \approx 166.3\,\%$.
\textbf{3.} From 2005 to 2010 ($26\,270$ to $26\,170$) and from 2015 to 2020 ($28\,390$ to
$28\,370$).
\textbf{4.} $(27\,700 - 22\,730) \div 22\,730 = 4\,970 \div 22\,730 \approx 21.9\,\%$.
\textbf{5.} $1.01^{25} \approx 1.282$, close to Brenland's factor $31\,470 \div 24\,680 \approx 1.275$;
$1.04^{25} \approx 2.666$, close to Kestria's factor $27\,700 \div 10\,400 \approx 2.663$.
\textbf{6.} Brenland $70 \div 1 = 70$ years; Kestria $70 \div 4 = 17.5$ years.
\textbf{7.} Brenland about $2025 + 70 = 2095$; Kestria about $2025 + 17.5$, that is around 2042.
\textbf{8.} It can copy methods and machines already used in richer countries, and where there is
little capital, each new machine adds a lot to output.
\textbf{9.} Brenland: $31\,470 \times 1.01^{10} \approx 34\,760$; Kestria: $27\,700 \times 1.04^{10}
\approx 41\,000$ dalens. Kestria's is higher.
\textbf{10.} 2030: Brenland $31\,470 \times 1.01^{5} \approx 33\,080$, Kestria $27\,700 \times
1.04^{5} \approx 33\,700$. Kestria is already ahead in 2030: it passes Brenland between 2025 and 2030.
\textbf{11.} 2029 ($n = 4$): Brenland $\approx 32\,748$, Kestria $\approx 32\,405$. 2030 ($n = 5$):
Brenland $\approx 33\,075$, Kestria $\approx 33\,701$. The first year is 2030.
\textbf{12.} For example: catching-up countries slow down as they get closer to the richest (the
breaks of the model), and a crisis or a war can cut \omterm{def:g7:households:income}{income} for years.
\end{solution}
